\section{The final bound}\label{sec:final-bill}

We keep the notation of Sections~\ref{sec:introduction} and
\ref{sec:notation}.
All error terms below are nonnegative.
The analytic inputs are \textup{(A1)}--\textup{(A5)} of
Table~\ref{tab:intro-inputs}.

\subsection{Addition of the comparisons}

Put $\delta_{\mathrm{err}}=10^{-4}$.
The constant of \eqref{eq:intro-DB} is
\begin{equation}\label{eq:bill-constant}
\begin{aligned}
 D_B={}&D_{\mathrm{II}}+\frac{5719}{10^6}
       +M_{5,\mathrm{internal}}+g_{\mathrm{RA,lo}}\\
     &-\delta_{\mathrm{err}}+
       \left(2-\frac{c_{K2,G0}}{4}\right)\overline H,
 \qquad \overline H=\frac{389253}{5000000}.
\end{aligned}
\end{equation}
Here $D_{\mathrm{II}}$ is Paper~II's original lower constant
\cite[(6.3a)]{PaperII}.
The six endpoint gains $g_6$ of \eqref{eq:retB-six-gain} are not used.
The two profile gains enter separately.
In particular,
\begin{equation}\label{eq:bill-components}
\begin{aligned}
 D_{\mathrm{II}}&=\frac{45849916338009170593}
                         {25000000000000000000},\\
 M_{5,\mathrm{internal}}
     &=-K_{\mathrm{hi}}\frac{6073714261}{10^9}-L_{G1},\\
 g_{\mathrm{RA,lo}}&=\frac{989598353}{10^{10}}=0.0989598353.
\end{aligned}
\end{equation}
The last line is the conservative rational floor for the integral
gain $\mathfrak g_{\mathrm{RA}}$. It is not a central estimate.
Section~\ref{sec:root-a} separates it from the integral gain.
The certified integral enclosure is
$[0.0989598353353,0.0989598353355]$.
We take $K_{\mathrm{hi}}$ and $L_{G1}$ as defined in
\cite[(A.10d), (4.84)]{PaperII}.
Their exact values are
\begin{equation}\label{eq:bill-M5-inputs}
\begin{aligned}
 K_{\mathrm{hi}}&=
 \frac{3989480791070578466390443009830933753748442947478110112223}
      {10^{58}},\\[2pt]
 L_{G1}&=-
 \frac{243950428265194656254206921939947762868871361874409552095097}
      {100433627766186892221372630771322662657637687111424552206336}.
\end{aligned}
\end{equation}
Thus $M_{5,\mathrm{internal}}=M_n/M_d$ exactly, where
\begin{equation}\label{eq:bill-M5-integers}
 M_n=\sum_{i=0}^5m_i10^{20i},\qquad
 M_d=\sum_{i=0}^5d_i10^{20i}.
\end{equation}
The coefficients are printed in Table~\ref{tab:bill-M5-integers}.
\begingroup
\small
\setlength{\tabcolsep}{8pt}
\renewcommand{\arraystretch}{1.25}
\begin{longtable}{@{}rrr@{}}
\caption{Exact shifted-profile gain}\label{tab:bill-M5-integers}\\
\hline
$i$ & $m_i$ & $d_i$\\
\hline
\endfirsthead
\hline
$i$ & $m_i$ & $d_i$\\
\hline
\endhead
0 & 66922148532305907989 & 0\\
1 & 34893223604885136721 & 0\\
2 & 3973350097562772509 & 0\\
3 & 87266878234046404033 & 48635364229120000000\\
4 & 28374567121249055735 & 73384187692692674921\\
5 & 3998 & 680564\\
\hline
\end{longtable}
\endgroup
The exact rational $c_{\mathrm{O1}}$ is defined in
\eqref{eq:l3-exact-certificate-constant} and
Table~\ref{tab:l3-certificate-integers}.
It records the optional sharper certificate only.
The main count uses $c_{K2,G0}$ in both comparisons.
The common functional comparison is \eqref{eq:l3-certificate}.
For comparison with the optional sharper coefficient, exact rational
arithmetic gives
\begin{equation}\label{eq:bill-coefficient-cost}
 0<c-c_{\mathrm{O1}}<4.081\cdot10^{-46},\qquad
 0<\frac{c-c_{\mathrm{O1}}}{4}\overline H
                  <7.943\cdot10^{-48}.
\end{equation}
This is a rational comparison only. It gives no bound for $2e^{-\gamma}U_0$.
The further cost of the common coefficient relative to \(c\) is exactly
\[
 \frac{c_{K2,G0}-c}{4}\overline H
 =0.000049489563023753677508253183146688036059673516909.
\]
The resulting main switched gain is exactly
\[
 \Delta D_{L3}=0.02263846472429815434060381043792394540367262995113415.
\]
Neither optional certificate is required for \eqref{eq:bill-constant}.
Each gain is proved on the original arithmetic source of its comparison.
The gains $5719/10^6$ and $M_{5,\mathrm{internal}}$ come from
Proposition~\ref{prop:retB-comparison}, through
Propositions~\ref{prop:app-h6-four-profiles} and~\ref{prop:app-m5-cap}.
The gain $g_{\mathrm{RA,lo}}$ comes from Proposition~\ref{prop:ra-b11}, and
$\Delta D_{L3}$ from \eqref{eq:l3-final}. Table~\ref{tab:bill-binding} lists the bindings.

Table~\ref{tab:bill-components} itemizes the constant.
Its decimal column gives approximations only.
The defining formulas determine the exact sum.
\begingroup
\small
\renewcommand{\arraystretch}{1.3}
\begin{longtable}{@{}p{.48\linewidth}p{.43\linewidth}@{}}
\caption{The six terms in $D_B$}\label{tab:bill-components}\\
\hline
Term & Value\\
\hline
\endfirsthead
\hline
Term & Value\\
\hline
\endhead
$D_{\mathrm{II}}$ & $1.83399665352036682372$\\
$5719/10^6$ & $0.005719$\\
$M_{5,\mathrm{internal}}$ & $0.005874949945026357\ldots$\\
$g_{\mathrm{RA,lo}}$ & $0.0989598353$\\
$-\delta_{\mathrm{err}}$ & $-0.0001$\\
$\Delta D_{L3}=(2-c_{K2,G0}/4)\overline H$
 & $0.022638464724298154\ldots$\\
\hline
$D_B$ & $1.967088903489691335\ldots$\\
\hline
\end{longtable}
\endgroup

Adding \eqref{eq:intro-comparisons} and
\eqref{eq:intro-exterior} cancels the sixth endpoint:
\begin{equation}\label{eq:bill-cancellation}
 \overline J_6-\overline J_6=0,\qquad
 2\overline H-\frac{c_{K2,G0}}{4}\overline H
 =\left(2-\frac{c_{K2,G0}}{4}\right)\overline H.
\end{equation}
Thus their conditional consequence is
\begin{equation}\label{eq:bill-added}
 \frac{4G_b(N)}{W(N)}
 \ge D_B+\delta_{\mathrm{err}}-c_{K2,G0}I(b)-\mathcal E_B(N;b,\sigma).
\end{equation}
This uses the complete $S_2$ of Paper~II.
The finite identities keep their original signs and multiplicities.
The fixed errors in Paper~II are subtracted once
\cite[(6.4)--(6.6)]{PaperII},
and not again in \eqref{eq:bill-constant}.

\subsection{Ten error allocations}

Use the exterior error from \eqref{eq:bill-K-error}.
For the table below, a conservative majorant is convenient:
\begin{equation}\label{eq:bill-K-error}
\begin{aligned}
 E_K(N;b,\sigma)
 &=A_K(b)C_{\mathrm{AA}}\sigma
      +\frac{3884000000000}{9999}\sigma^9+r_K(N;b,\sigma)\\
 &\le 4A_K(b)C_{\mathrm{AA}}\sigma
      +\frac{3884000000000}{9999}\sigma^9
      +r_K(N;b,\sigma),\\
 A_K(b)&=2e^{-\gamma}I(b)W_{\mathrm{AA},K},\qquad
 P_K=\frac{3884000000000}{9999},\\
 r_K(N;b,\sigma)&=o_{b,\sigma}(1).
\end{aligned}
\end{equation}
Here $W_{\mathrm{AA},K}$ is the exterior remainder-weight bound.
It belongs to the source in Section~\ref{sec:k2}.
The extra factor four only enlarges this nonnegative error.
It leaves the principal constant unchanged.
Similarly, the switched comparison requires
\begin{equation}\label{eq:bill-sw-error}
\begin{aligned}
 E_{\mathrm{sw}}(N;\sigma)
 &\le \frac{650000000}{1111}C_E\sigma
       +\frac{48544888063158}{27775}\sigma^9
       +V_{\mathrm{sw}}(N;\sigma),\\
 V_{\mathrm{sw}}(N;\sigma)&=o_\sigma(1).
\end{aligned}
\end{equation}
Here $V_{\mathrm{sw}}$ is the positive part in
\eqref{eq:l3-remainder}.
The conditional limit is \eqref{eq:l3-remainder-limit}.
These two limits concern the actual sums stated there.
Both majorants are the error terms of \eqref{eq:k2-final-fixed-sigma} under
\textup{(A3)} and of \eqref{eq:l3-final} under \textup{(A4)}, with
$C_{\mathrm{sw}}$ from \eqref{eq:l3-price}. By \textup{(A5)},
$C_{\mathrm{AA}}$ and $C_E$ do not depend on $\sigma$. The exterior weight
$W_{\mathrm{AA},K}$ is $W_{\mathrm{AA}}$ of \eqref{eq:k2-aa-mass}:
\begin{equation}\label{eq:bill-AK-positive}
\begin{gathered}
 A_K(b)=2e^{-\gamma}I(b)W_{\mathrm{AA},K},\qquad
 W_{\mathrm{AA},K}=500\Bigl(1+H_{\log}+\frac{H_{\log}^2}2
                    +\frac{H_{\log}^3}6\Bigr),\\
 H_{\log}=\log\frac{\mu_*}{u}=\log\frac{3125}{32}>0.
\end{gathered}
\end{equation}
Since $I(b)>0$, this gives $0<A_K(b)<\infty$. In particular
$A_K(b)C_{\mathrm{AA}}\sigma\le4A_K(b)C_{\mathrm{AA}}\sigma$.

The retained and Root--A comparisons require
\begin{equation}\label{eq:bill-eventual-errors}
 E'(N;\sigma)=o_\sigma(1),\qquad
 E_{\mathrm{RA}}(N;\sigma)=o_\sigma(1).
\end{equation}
For Root--A, use $E_{\mathrm{RA}}(N;\sigma)$
from \eqref{eq:ra-b11-floor}.
The first includes every conversion and the saved-array contribution
$\sum_{j\in\mathcal H}|\Lambda^0_{j,N}|$ in
\eqref{eq:intro-retained-error}; that contribution is not charged again.
The pointwise error may use the coarser bound
\begin{equation}\label{eq:bill-pointwise-error}
 E_{\mathrm{exc}}(N)=\frac{4E_{\mathrm{pt}}(N)}{W(N)}
 \le\frac{2197N^{1-\alpha}}{W(N)}=o(1).
\end{equation}
Here $E_{\mathrm{pt}}$ is defined in
\eqref{eq:source-pointwise-error}. The explicit pointwise argument in
Section~\ref{sec:source-decomposition} gives $E_{\mathrm{pt}}=0$
for $N\ge N_{\mathrm{pt}}$. The coarser majorant retains the same error bound.
The inequality follows from \eqref{eq:source-first-lower}, using both
the pointwise bound and the exceptional-set cardinality.
Since $C(N)\ge C_2>0$ and $\alpha>0$, the coarser majorant tends to zero:
$N^{-\alpha}(\log N)^2\to0$.
No sharper exceptional coefficient is needed.
Proposition~\ref{prop:retB-comparison} gives $E'(N;\sigma)\to0$, and
Proposition~\ref{prop:ra-b11} gives $E_{\mathrm{RA}}(N;\sigma)\to0$ under
\textup{(A2)}, each at fixed $\sigma$.
Lemma~\ref{lem:source-pointwise} gives $E_{\mathrm{exc}}(N)=0$ for
$N\ge N_{\mathrm{pt}}$.

Fix $b$, all absolute constants, and the source geometry first.
Choose positive absolute constants $C_E^{\mathrm{L3}}$ and
$C_E^{\mathrm{II}}$ for the Section~\ref{sec:l3} and Paper~II sieve bounds. Define
\begin{equation}\label{eq:bill-CE-definition}
 C_E:=\max\{C_E^{\mathrm{L3}},C_E^{\mathrm{II}}\}.
\end{equation}
By \textup{(A5)}, these constants are independent of $\sigma$ and $N$.
Here $C_{\mathrm{AA}}=4C_{\mathrm{AA37}}$ by \eqref{eq:k2-aa-constant}. The
constant $C_E$ dominates the absolute constant of
\eqref{eq:classical-structured-sieve} used in \eqref{eq:l3-sieve-error}.
The underlying sieve bounds come from Paper~I, Theorem~2.6, (2.31)
\cite[Theorem~2.6, (2.31)]{PaperI}.
Choose one rational $\sigma_*>0$ satisfying
\begin{equation}\label{eq:bill-sigma}
\begin{aligned}
 0<\sigma_*&\le 10^{-10},\\
 \sigma_*&\le\frac{713332}{625\cdot10^{21}C_E},\\
 \sigma_*&\le\frac{1}{400000A_K(b)C_{\mathrm{AA}}},\\
 \sigma_*&\le\left(\frac{9999}{388400000000000000}\right)^{1/9},\\
 \sigma_*&\le\frac{1111}{65000000000000C_E},\\
 \sigma_*&\le\left(\frac{27775}{4854488806315800000}\right)^{1/9},\\
 \tau_*&=(\sigma_*)^9.
\end{aligned}
\end{equation}
By \eqref{eq:bill-AK-positive} and \textup{(A5)}, every right side in
\eqref{eq:bill-sigma} is positive and independent of $\sigma$.
Hence such a rational $\sigma_*$ exists.
The second restriction preserves Paper~II's fixed error bound.
It is exactly
\begin{equation}\label{eq:bill-paperII-AA}
 \sigma_*\le\frac{4g_{\mathrm{AA}}}{C_EW_*},\qquad
 g_{\mathrm{AA}}=\frac{178333}{6250000000},\qquad W_*=10^{14}
 \quad\text{\cite[(1.26), (1.27), (6.6)]{PaperII}}.
\end{equation}
The other four restrictions control new fixed errors.
Freeze every finite partition and logarithmic precision next.
Only then choose the common threshold.

The ten allocations are printed in Table~\ref{tab:bill}.
The two zero rows are unused capacity.
\begingroup
\small
\renewcommand{\arraystretch}{1.65}
\setlength{\extrarowheight}{2pt}
\begin{longtable}{@{}r p{.60\linewidth} p{.25\linewidth}@{}}
\caption{The error bound $\delta_{\mathrm{err}}=10^{-4}$}
\label{tab:bill}\\
\hline
 & Error term & Allocation\\
\hline
\endfirsthead
\hline
 & Error term & Allocation\\
\hline
\endhead
1 & $4A_K(b)C_{\mathrm{AA}}\sigma_*$ & $\delta_{\mathrm{err}}/10$\\
2 & $\tfrac{3884000000000}{9999}\tau_*$ & $\delta_{\mathrm{err}}/10$\\
3 & $\tfrac{650000000}{1111}C_E\sigma_*$ & $\delta_{\mathrm{err}}/10$\\
4 & $\tfrac{48544888063158}{27775}\tau_*$ & $\delta_{\mathrm{err}}/10$\\
5 & $E'(N;\sigma_*)+V_{\mathrm{sw}}(N;\sigma_*)$
    & $\tfrac{\delta_{\mathrm{err}}}{20}+\tfrac{\delta_{\mathrm{err}}}{20}$\\
6 & $0$ & $\delta_{\mathrm{err}}/10$\\
7 & $E_{\mathrm{RA}}(N;\sigma_*)$ & $\delta_{\mathrm{err}}/10$\\
8 & $r_K(N;b,\sigma_*)$ & $\delta_{\mathrm{err}}/10$\\
9 & $E_{\mathrm{exc}}(N)$ & $\delta_{\mathrm{err}}/10$\\
10 & $0$ & $\delta_{\mathrm{err}}/10$\\
\hline
\end{longtable}
\endgroup
Its capacities sum exactly:
\begin{equation}\label{eq:bill-sum}
 4\frac{\delta_{\mathrm{err}}}{10}
 +\left(\frac{\delta_{\mathrm{err}}}{20}
       +\frac{\delta_{\mathrm{err}}}{20}\right)
 +5\frac{\delta_{\mathrm{err}}}{10}
 =\delta_{\mathrm{err}}.
\end{equation}
The first four bounds follow from \eqref{eq:bill-sigma}.
The five remaining nonzero bounds follow from the stated limits.
Each limit is taken at fixed $\sigma_*$.

\begin{proposition}[Error budget]\label{prop:bill-budget}
Assume \textup{(A1)}--\textup{(A5)}, and let $\sigma_*$ be rational with
\eqref{eq:bill-sigma}. Then every even $N\ge N_0^*(\sigma_*)$ satisfies
\begin{equation}\label{eq:bill-budget}
 \mathcal E_B(N;b,\sigma_*)\le\frac8{10}\delta_{\mathrm{err}}
 <\delta_{\mathrm{err}}.
\end{equation}
\end{proposition}

\begin{proof}
Rows 1 and 3. The third and fifth lines of \eqref{eq:bill-sigma} give
\[
 4A_K(b)C_{\mathrm{AA}}\sigma_*\le\frac4{400000}=\frac{\delta_{\mathrm{err}}}{10},
 \qquad
 \frac{650000000}{1111}C_E\sigma_*\le\frac{650000000}{65000000000000}
 =\frac{\delta_{\mathrm{err}}}{10}.
\]
Rows 2 and 4. The first line gives $\tau_*\le10^{-90}$. This is below
$9999/388400000000000000$ and $27775/4854488806315800000$, so the fourth and
sixth lines hold, and
\[
 \frac{3884000000000}{9999}\tau_*\le\frac{\delta_{\mathrm{err}}}{10},\qquad
 \frac{48544888063158}{27775}\tau_*\le\frac{\delta_{\mathrm{err}}}{10}.
\]
Rows 5, 7 and 8. At the fixed $\sigma_*$, $E'\to0$ by \textup{(A1)},
$V_{\mathrm{sw}}\to0$ by \eqref{eq:l3-remainder-limit} under \textup{(A4)},
$E_{\mathrm{RA}}\to0$ by \textup{(A2)}, and $r_K\to0$ by \textup{(A3)}.
By \eqref{eq:not-little-o}, the thresholds
$T_{E'},T_V,T_{E_{\mathrm{RA}}},T_r,T_{E_{\mathrm{exc}}}$ of
\eqref{eq:not-common-threshold} exist,
and above them $E'\le\delta_{\mathrm{err}}/20$,
$V_{\mathrm{sw}}\le\delta_{\mathrm{err}}/20$,
$E_{\mathrm{RA}}\le\delta_{\mathrm{err}}/10$ and
$r_K\le\delta_{\mathrm{err}}/10$.
Row 9. Lemma~\ref{lem:source-pointwise} gives $E_{\mathrm{exc}}(N)=0$ for
$N\ge N_{\mathrm{pt}}$, so $T_{E_{\mathrm{exc}}}\le N_{\mathrm{pt}}$.
Rows 6 and 10 are zero.
Every threshold used is bounded by the common maximum in
\eqref{eq:not-common-threshold}. Also $N_{\mathrm{pt}}<N_{\mathrm{tr}}$
by Lemma~\ref{lem:bill-transfer}. Hence all of them hold for
$N\ge N_0^*(\sigma_*)$. By \eqref{eq:intro-error-identifications},
\eqref{eq:bill-K-error} and \eqref{eq:bill-sw-error},
\[
\begin{aligned}
 \mathcal E_B
 \le{}&\bigl(E'+V_{\mathrm{sw}}\bigr)+E_{\mathrm{RA}}
   +\Bigl(\frac{650000000}{1111}C_E\sigma_*+\frac{48544888063158}{27775}\tau_*\Bigr)\\
   &+\Bigl(4A_KC_{\mathrm{AA}}\sigma_*+\frac{3884000000000}{9999}\tau_*+r_K\Bigr)
   +E_{\mathrm{exc}}.
\end{aligned}
\]
The bounds above give at most
$\frac{\delta_{\mathrm{err}}}{10}+\frac{\delta_{\mathrm{err}}}{10}
+\frac{2\delta_{\mathrm{err}}}{10}+\frac{3\delta_{\mathrm{err}}}{10}+0$,
which is below the nonzero allocations
$\frac{8\delta_{\mathrm{err}}}{10}$ of Table~\ref{tab:bill}.
\end{proof}

Take the named source, application and error thresholds listed in
Section~\ref{sec:notation}. The Section~\ref{sec:k2} definitions of
\(T_K\) and \(T_r\) already include the main-mass threshold
\(N_{\mathrm{mass}}(b,\sigma_*)\) and the AA threshold
\(N_{\mathrm{AA}}(b,\sigma_*)\). Thus the common maximum in \(\mathcal T_*\)
below includes both requirements at the fixed \(b,\sigma_*\).
The common threshold is the single definition
\begin{equation}\label{eq:bill-threshold}
 N_0^*(\sigma_*)\quad\text{as defined in \eqref{eq:not-common-threshold}}.
\end{equation}
Its existence depends on all those estimates. The floor alone supplies
none of them. By its definition and Lemma~\ref{lem:bill-transfer},
$N_0^*(\sigma_*)\ge\max\{N_{\mathrm{tr}},N_6\}>N_{\mathrm{pt}}$.
The Section~\ref{sec:notation} threshold definitions include $N_{\mathrm{pt}}$ through
$T_{\mathrm{src}}$ and $N_6$ through $T_{\mathrm{RA}}$. The bound $N_{\mathrm{tr}}$
enters directly in the common maximum. Lemma~\ref{lem:source-pointwise}
holds for $N\ge N_{\mathrm{pt}}$ and Lemma~\ref{lem:source-upper-sources}
for $N\ge N_6$.
Thus the order is
\begin{equation}\label{eq:bill-order}
 \text{constants}\ \longrightarrow\ \sigma_*
 \ \longrightarrow\ \text{finite choices}
 \ \longrightarrow\ N_0^*(\sigma_*)\ \longrightarrow\ N.
\end{equation}
No parameter $\sigma_*(N)$ is used.

Proposition~\ref{prop:bill-budget} gives
\begin{equation}\label{eq:bill-error-bound}
 \mathcal E_B(N;b,\sigma_*)\le\delta_{\mathrm{err}}
 \qquad(N\ge N_0^*(\sigma_*),\ N\in2\mathbb Z).
\end{equation}
Combining \eqref{eq:bill-added} and \eqref{eq:bill-error-bound} gives
\begin{equation}\label{eq:bill-positive-count}
 \boxed{\displaystyle
 \frac{4G_b(N)}{W(N)}\ge D_B-c_{K2,G0}I(b).}
\end{equation}
It holds for every even $N\ge N_0^*(\sigma_*)$ under
\textup{(A1)}--\textup{(A5)}.

\subsection{The numerical margin}

All constants in \eqref{eq:bill-constant} have exact definitions.
The exterior price uses the G0 coefficient in
\eqref{eq:k2-certificate}.
The same coefficient prices L3. No $G_*,U_0$ sharper domination
is used in this arithmetic or in the main-count hypotheses.
The constants have the following outward enclosures:
\begin{equation}\label{eq:bill-numeric-enclosures}
\begin{gathered}
 1.967088903489691335582703173494<D_B,\\
 D_B<1.967088903489691335582703173495,\\[2pt]
 1.966859949178424046530804965343<c_{K2,G0}I(b),\\
 c_{K2,G0}I(b)<1.966859949178424046530804965344,\\[2pt]
 0.000228954311267289051898208151<D_B-c_{K2,G0}I(b),\\
 D_B-c_{K2,G0}I(b)<0.000228954311267289051898208152.
\end{gathered}
\end{equation}
Here
\begin{equation}\label{eq:bill-price-definition}
 c_{K2,G0}=\frac{6836825163875517756286846321650754372931094689}{10^{45}},
 \qquad I(b)=\log\frac{4000007}{2999993}.
\end{equation}
The coefficient is defined in \eqref{eq:k2-certificate}.
The final comparison is \eqref{eq:k2-final-fixed-sigma}.
Hence
\begin{equation}\label{eq:bill-margin-simple}
 D_B-c_{K2,G0}I(b)>0.0002289>\frac1{5000}>0.
\end{equation}
The table below identifies each exact constant by its defining display and
lists the comparisons that use it unchanged.
\begingroup
\small
\renewcommand{\arraystretch}{1.3}
\begin{longtable}{@{}p{.27\linewidth}p{.31\linewidth}p{.34\linewidth}@{}}
\caption{Constants and the comparisons that use them}
\label{tab:bill-binding}\\
\hline
Constant & Definition & Comparison\\
\hline
\endfirsthead
\hline
Constant & Definition & Comparison\\
\hline
\endhead
$D_{\mathrm{II}}$ & \eqref{eq:intro-components}, \cite[(6.3a)]{PaperII}
 & Proposition~\ref{prop:retB-comparison}, by \eqref{eq:retB-paperII-floor}\\
$5719/10^6$ & \eqref{eq:retB-H6-finite}
 & Proposition~\ref{prop:retB-comparison}, by
   Proposition~\ref{prop:app-h6-four-profiles}\\
$M_{5,\mathrm{internal}}$, $K_{\mathrm{hi}}$, $L_{G1}$
 & \eqref{eq:bill-M5-inputs}, Table~\ref{tab:bill-M5-integers}
 & Proposition~\ref{prop:retB-comparison}, by
   Proposition~\ref{prop:app-m5-cap}\\
$g_{\mathrm{RA,lo}}$ & \eqref{eq:ra-rational-floor}
 & Proposition~\ref{prop:ra-b11}, \eqref{eq:ra-b11-floor}\\
$\overline H$, $c_{K2,G0}$ & \eqref{eq:bill-constant}, \eqref{eq:k2-certificate}
 & \eqref{eq:l3-final}, with \eqref{eq:l3-price}\\
$c_{K2,G0}$, $I(b)$ & \eqref{eq:bill-price-definition}
 & Proposition~\ref{prop:k2-comparison}, \eqref{eq:k2-final-fixed-sigma}\\
$A_K(b)$, $C_{\mathrm{AA}}$, $P_K$
 & \eqref{eq:bill-AK-positive}, \eqref{eq:k2-aa-constant},
   \eqref{eq:k2-final-fixed-sigma}
 & Table~\ref{tab:bill}, rows 1 and 2\\
$C_{\mathrm{sw}}$, $P_{\mathrm{sw}}$ & \eqref{eq:l3-price}
 & Table~\ref{tab:bill}, rows 3 and 4\\
$\delta_{\mathrm{err}}$ & \eqref{eq:bill-constant}
 & Proposition~\ref{prop:bill-budget}\\
\hline
\end{longtable}
\endgroup
No other constant enters $D_B$ or the error budget.

The logarithm enclosure uses rational arithmetic.
Put $t_{\log}=500007/3500000$ and
\begin{equation}\label{eq:bill-log-enclosure}
 S_m=2\sum_{k=0}^{m}\frac{t_{\log}^{2k+1}}{2k+1},
 \qquad
 S_m< I(b)<S_m+
       \frac{2t_{\log}^{2m+3}}{(2m+3)(1-t_{\log}^2)}.
\end{equation}
Here $I(b)=2\operatorname{arctanh}t_{\log}$ and $0<t_{\log}<1$.
Bounding each remaining denominator by $2m+3$ proves the upper bound.
Taking $m=24$ gives \eqref{eq:bill-numeric-enclosures}.
Thus no rounded logarithm enters the margin test.

For reference, put $\zeta=10^{-6}$ and $\mathfrak E_*=\exp(D_B/c_{K2,G0})$.
The buffered critical exponent is
\begin{equation}\label{eq:bill-critical-exponent}
 a_{\mathrm{buf}}^*
 =\frac{\mathfrak E_*+1}{(1-\zeta)\mathfrak E_*-\zeta}.
\end{equation}
To derive it, write $b_*=1/(1+\mathfrak E_*)+\zeta$.
The equation $a(1-b_*)=1$ gives
\eqref{eq:bill-critical-exponent}.
For the present constants,
\begin{equation}\label{eq:bill-critical-enclosure}
 1.749974884137892433433207484622
 <a_{\mathrm{buf}}^*
 <1.749974884137892433433207484623<\frac74.
\end{equation}
This enclosure also uses rational bounds.
For $x=D_B/c_{K2,G0}$ and $m=36$, put
\begin{equation}\label{eq:bill-exp-enclosure}
 T_m=\sum_{k=0}^{m}\frac{x^k}{k!},\qquad
 T_m<e^x<T_m+
 \frac{x^{m+1}}{(m+1)!}\frac{1}{1-x/(m+2)}.
\end{equation}
Here $0<x<1$.
The tail ratio is at most $x/(m+2)$.
The function in \eqref{eq:bill-critical-exponent} decreases with $\mathfrak E_*$.
Substitution of the two rational bounds proves
\eqref{eq:bill-critical-enclosure}.

The endpoint in \eqref{eq:bill-critical-enclosure} is equivalent to the
positive scalar margin at this fixed $b$. It supplies no theorem at a
smaller $a$ without new comparisons at the corresponding cutoff.
If the optional Lemma~\ref{lem:retB-six-endpoints} is used, $D_B$ increases by
$g_6=0.0020549454896034927680\ldots$ of \eqref{eq:retB-six-gain}. The margin
then becomes $0.0022838998008707818199\ldots$, and $a_{\mathrm{buf}}^*$ becomes
$1.7497494980881623276859\ldots$. Theorem~\ref{thm:main} does not use this.

For an optional arithmetic comparison, define
$D_{B,c}=D_B+(c_{K2,G0}-c)\overline H/4$.
Using $c$ in both K2 and L3 would require two further inputs.
This paragraph concerns the optional scalar comparison at the same fixed cutoff;
it supplies no new cutoff-dependent estimate. The first is
the separate sharper functional domination $G_*\le U_0$ from
\hyperref[eq:app-g24-inclusion]{equation~\ref*{eq:app-g24-inclusion}}. The second is a certified
$2e^{-\gamma}U_0\le c$. It does not follow from the coarse
interval in \hyperref[eq:app-g24-numbers]{equation~\ref*{eq:app-g24-numbers}}. The main proof uses
neither.
Its scalar endpoint is
\begin{equation}\label{eq:bill-active-critical-exponent}
 E_c=\exp(D_{B,c}/c),\qquad
 a^*_{\mathrm{buf},c}
 =\frac{E_c+1}{(1-\zeta)E_c-\zeta}.
\end{equation}
The same rational Taylor enclosure gives
\begin{equation}\label{eq:bill-active-critical-enclosure}
 1.749889173139093190128834678
 <a^*_{\mathrm{buf},c}
 <1.749889173139093190128834679<\frac74.
\end{equation}
Thus, on the optional sharper branch, $a^*_{\mathrm{buf},c}$ is
$1.7498891731390\ldots$. The main theorem uses the explicitly certified
G0 coefficient defined and bounded in \eqref{eq:k2-certificate}; the
display above is only the optional exponent endpoint.

\subsection{The transfer}

\begin{lemma}[Transfer]\label{lem:bill-transfer}
The parameters \eqref{eq:not-transfer-parameters} satisfy
\begin{equation}\label{eq:bill-transfer-parameters}
 \kappa_{\mathrm{tr}}=\frac34(1-b)-b=\frac7{4000000},\qquad
 N_{\mathrm{tr}}=10^{12000000/7}>10^{12}+1=N_{\mathrm{pt}}.
\end{equation}
For even $N\ge N_{\mathrm{tr}}$, every $n=rq\in\mathcal A_N$ with primes
$r\le q$ and $r<N^b$ satisfies
\begin{equation}\label{eq:bill-transfer}
 q>\varepsilon N^{1-b},\qquad
 q^{3/4}>\varepsilon^{3/4}N^{b+\kappa_{\mathrm{tr}}}\ge N^b>r.
\end{equation}
\end{lemma}

\begin{proof}
Since $a=7/4$, $a(1-b)-1=\frac34(1-b)-b$, and $b=2999993/7000000$ gives
$7/4000000$. Also $\varepsilon^{-(a-1)/\kappa_{\mathrm{tr}}}
=10^{4\cdot3000000/7}$, and $12000000/7>13$.
For $n=rq\in\mathcal A_N$ we have $n>\varepsilon N$, so
$q=n/r>\varepsilon N^{1-b}$ and
$q^{3/4}>\varepsilon^{3/4}N^{3(1-b)/4}=\varepsilon^{3/4}N^{b+\kappa_{\mathrm{tr}}}$.
Finally $N^{\kappa_{\mathrm{tr}}}\ge\varepsilon^{-3/4}$ exactly when
$N\ge N_{\mathrm{tr}}$.
\end{proof}

\section*{Proof of Theorem 1}

Use the established inputs \textup{(A1)}--\textup{(A5)} and fix a rational
$\sigma_*$ satisfying \eqref{eq:bill-sigma}. Let
$N\ge N_0^*(\sigma_*)$ be even. Then
$N\ge\max\{N_{\mathrm{tr}},N_6\}>N_{\mathrm{pt}}$.
Lemma~\ref{lem:source-pointwise} gives the first line of
\eqref{eq:intro-comparisons}. At $\sigma=\sigma_*$, the established
comparisons \textup{(A1)}, \textup{(A2)} and \textup{(A4)} give its other
three lines, and \textup{(A3)} gives \eqref{eq:intro-exterior}.
The zero-gain combination
Lemma~\ref{lem:source-trackb-combination} adds these four comparisons to
\eqref{eq:bill-added}. Proposition~\ref{prop:bill-budget} gives
\eqref{eq:bill-error-bound}, hence \eqref{eq:bill-positive-count}.
Equation~\eqref{eq:bill-margin-simple} then gives $G_b(N)>0$, since
$W(N)>0$ by \eqref{eq:not-normalization}.
Hence the set in \eqref{eq:intro-count} is nonempty.

Choose one counted integer $n=N-p$.
If $n$ is prime, take $q=n$ and $r=1$.
Otherwise $n=rq$ with primes $r\le q$ and $r<N^b$, and
Lemma~\ref{lem:bill-transfer} gives $r<q^{3/4}$.
Therefore $N=p+rq$ and $r\le q^{3/4}$. $\square$
